% Seção DevOps — NitrusLeaf Mobile
% Incluir no documento principal: % Seção DevOps — NitrusLeaf Mobile
% Incluir no documento principal: % Seção DevOps — NitrusLeaf Mobile
% Incluir no documento principal: % Seção DevOps — NitrusLeaf Mobile
% Incluir no documento principal: \input{docs/devops-nitrusleaf.tex}
% Requer no preâmbulo: \usepackage{tikz}, \usetikzlibrary{positioning,arrows.meta}
%                       \usepackage{float}, \usepackage{enumitem}

\section{DevOps: Infraestrutura, Automação e Entrega Contínua}

DevOps \cite{wang2023} é uma cultura e conjunto de práticas que integra desenvolvimento de software (Dev) e operações de TI (Ops), focando em automação, colaboração e entrega contínua. No \textit{NitrusLeaf Mobile}, o repositório \texttt{Vitalliz/nitrusleaf-mobile} adota Git, GitHub Actions, backend Supabase e aplicativo Expo (React Native), com validação automática a cada alteração enviada ao repositório.

\subsection{Pipeline CI/CD}

O pipeline de integração e entrega contínua executa três \textit{jobs} no GitHub Actions (\texttt{Build}, \texttt{Test} e \texttt{Deploy}) a cada \textit{push} em qualquer branch ou em \textit{pull request} direcionado à \texttt{main}. A Figura~\ref{fig:cicd} resume o fluxo; a etapa \textbf{Monitor} representa a observação do resultado (painel do GitHub Actions e Supabase Dashboard) e o retorno de correções ao código --- ciclo de melhoria contínua.

\begin{figure}[H]
\centering
\begin{tikzpicture}[
    node distance=0.9cm and 1.1cm,
    stage/.style={
        rectangle, draw, rounded corners=3pt,
        fill=secondarycolor!20,
        text width=2.15cm, align=center,
        font=\small, minimum height=1.05cm
    },
    note/.style={font=\scriptsize, text=black!65, align=center},
    arr/.style={-{Stealth[length=2.2mm]}, thick}
]
% Linha principal
\node[stage] (code) {\textbf{Code}\\[-1pt]\scriptsize Git push / PR};
\node[stage, right=of code] (build) {\textbf{Build}\\[-1pt]\scriptsize \texttt{npm ci}};
\node[stage, right=of build] (test) {\textbf{Test}\\[-1pt]\scriptsize tipos + Jest};
\node[stage, right=of test] (deploy) {\textbf{Deploy}\\[-1pt]\scriptsize branch \texttt{main}};
\node[stage, right=of deploy] (monitor) {\textbf{Monitor}\\[-1pt]\scriptsize painéis};

\draw[arr] (code) -- (build);
\draw[arr] (build) -- (test);
\draw[arr] (test) -- (deploy);
\draw[arr] (deploy) -- (monitor);

% Detalhes abaixo (o que cada job executa de fato)
\node[note, below=0.55cm of build] (buildnote) {+\ \texttt{expo-doctor}\ (não bloqueia)};
\node[note, below=0.55cm of test] (testnote) {\texttt{typecheck},\ \texttt{lint:ci}};
\node[note, below=0.55cm of deploy] (deploynote) {integração após CI verde};

% Feedback Monitor -> Code
\draw[arr, dashed] (monitor.south) to[bend right=55]
    node[below, note, yshift=-0.15cm] {correções e novo commit} (code.south);

\end{tikzpicture}
\caption{Pipeline CI/CD implementado no repositório (\texttt{.github/workflows/ci.yml})}
\label{fig:cicd}
\end{figure}

\subsubsection{Correspondência entre etapas e jobs}

\begin{table}[H]
\centering
\caption{Etapas do pipeline e implementação no GitHub Actions}
\label{tab:cicd-jobs}
\small
\begin{tabular}{|l|l|p{7.2cm}|}
\hline
\textbf{Etapa} & \textbf{Job} & \textbf{Comandos / comportamento real} \\
\hline
Code & --- & \textit{Push} ou PR; código em \texttt{nitrusleaf-mobile/} \\
\hline
Build & \texttt{build} & \texttt{npm ci}; \texttt{npx expo-doctor} (\texttt{continue-on-error}) \\
\hline
Test & \texttt{test} & \texttt{npm run typecheck}; \texttt{npm test -- --ci}; \texttt{npm run lint:ci} \\
\hline
Deploy & \texttt{deploy} & Somente \textit{push} na \texttt{main}; confirma integração (EAS/Supabase opcionais) \\
\hline
Monitor & --- & GitHub Actions (status do workflow); Supabase Dashboard (API, Auth, BD) \\
\hline
\end{tabular}
\end{table}

O ambiente de execução utiliza \textbf{Node.js 22} no runner e a variável \texttt{FORCE\_JAVASCRIPT\_ACTIONS\_TO\_NODE24} para compatibilidade das ações oficiais do GitHub.

\subsection{Exemplo de Configuração CI/CD}

Trecho conceitual equivalente ao arquivo \texttt{.github/workflows/ci.yml}:

\begin{verbatim}
name: CI
on: [push, pull_request]

jobs:
  build:
    runs-on: ubuntu-latest
  steps:
    - uses: actions/checkout@v4
    - uses: actions/setup-node@v4
      with:
        node-version: "22"
    - run: npm ci
    - run: npx expo-doctor
      continue-on-error: true

  test:
    needs: build
  steps:
    - run: npm run typecheck
    - run: npm test -- --ci
    - run: npm run lint:ci

  deploy:
    needs: test
    if: branch == main
  steps:
    - run: echo "CI verde na main"
\end{verbatim}

\textbf{Comandos locais:}
\begin{itemize}[leftmargin=*, nosep]
    \item \texttt{npm run validate} --- verificação de tipos (\texttt{tsconfig.ci.json}) + testes Jest;
    \item \texttt{npm run lint:ci} --- ESLint na camada core (\texttt{utils}, \texttt{repositories}, \texttt{services}, \texttt{contexts});
    \item \texttt{npm run lint} --- análise completa do projeto (inclui telas; pode reportar avisos legados).
\end{itemize}

\subsection{Infraestrutura como Código}

\begin{table}[H]
\centering
\caption{Ferramentas DevOps utilizadas no NitrusLeaf Mobile}
\label{tab:devops-tools}
\begin{tabular}{|l|l|p{6cm}|}
\hline
\textbf{Categoria} & \textbf{Ferramenta} & \textbf{Finalidade} \\
\hline
Controle de Versão & Git + GitHub & \texttt{main} e branches da equipe via PR \\
CI/CD & GitHub Actions & Jobs \texttt{build}, \texttt{test} e \texttt{deploy} \\
Backend (BaaS) & Supabase & PostgreSQL, Auth JWT e API REST \\
App Mobile & Expo SDK 54 + React Native & iOS, Android e Web \\
Persistência local & \texttt{expo-sqlite} & SQLite no dispositivo (offline) \\
Testes & Jest + jest-expo & Testes em \texttt{src/utils/\_\_tests\_\_} \\
Qualidade & ESLint (\texttt{lint:ci}) & Validação automática da camada core \\
Configuração & \texttt{.env} / \texttt{.env.example} & Chaves Supabase (não versionadas) \\
\hline
\end{tabular}
\end{table}

\subsection{Infraestrutura do Sistema}

\subsubsection{Componentes por ambiente}

\begin{table}[H]
\centering
\caption{Componentes de infraestrutura}
\label{tab:infrastructure}
\begin{tabular}{|l|p{10cm}|}
\hline
\textbf{Componente} & \textbf{Descrição} \\
\hline
\textbf{Desenvolvimento local} &
\begin{itemize}[leftmargin=*, nosep]
    \item Node.js $\geq$ 20 (CI usa 22)
    \item Variáveis em \texttt{.env} conforme \texttt{.env.example}
    \item Branches dos membros integradas à \texttt{main} via pull request com CI verde
\end{itemize} \\
\hline
\textbf{Nuvem (Supabase)} &
\begin{itemize}[leftmargin=*, nosep]
    \item PostgreSQL gerenciado, autenticação e API REST
    \item HTTPS entre app e backend
    \item Painel para logs, usuários e consulta ao banco (etapa \textbf{Monitor})
\end{itemize} \\
\hline
\textbf{Dispositivos móveis} &
\begin{itemize}[leftmargin=*, nosep]
    \item Binário/servidor de desenvolvimento via Expo
    \item \texttt{expo-sqlite} para cache local
    \item Camada \texttt{src/repositories} para acesso local e remoto
\end{itemize} \\
\hline
\textbf{Segurança} &
\begin{itemize}[leftmargin=*, nosep]
    \item TLS nas requisições à API
    \item Sessão JWT via Supabase Auth
    \item Chave \textit{anon} no cliente; \texttt{.env} ignorado pelo Git
\end{itemize} \\
\hline
\end{tabular}
\end{table}

\subsubsection{Fluxo de dados}

\begin{itemize}
    \item \textbf{Usuário mobile:} autentica no Supabase Auth e consome dados via API.
    \item \textbf{Banco remoto:} PostgreSQL no Supabase.
    \item \textbf{Cache local:} SQLite (\texttt{expo-sqlite}) no aparelho.
    \item \textbf{CI/CD:} cada \textit{push} dispara build e testes; merge na \texttt{main} após pipeline bem-sucedido.
    \item \textbf{Monitoramento:} status dos workflows no GitHub; métricas e saúde da API no Supabase Dashboard.
\end{itemize}

\subsection{Arquitetura de Deploy}

A Figura~\ref{fig:deployment} separa validação do código (CI), execução do app e backend em nuvem. O GitHub Actions não hospeda o aplicativo: ele \textit{valida} o repositório antes da integração na \texttt{main}.

\begin{figure}[H]
\centering
\begin{tikzpicture}[
    node distance=1.5cm and 1.2cm,
    component/.style={
        rectangle, draw, rounded corners=3pt,
        fill=secondarycolor!20,
        text width=2.6cm, align=center,
        font=\small, minimum height=1.1cm
    },
    arr/.style={-{Stealth[length=2mm]}, thick},
    lbl/.style={font=\scriptsize, midway}
]
% Coluna CI
\node[component] (gh) {GitHub\\Repositório};
\node[component, below=of gh] (ci) {GitHub\\Actions};

% Coluna runtime
\node[component, right=3.2cm of gh] (app) {App Expo\\(dispositivo)};
\node[component, below=of app] (sqlite) {SQLite\\(\texttt{expo-sqlite})};

% Coluna backend
\node[component, right=3.2cm of app] (api) {Supabase\\Auth + API};
\node[component, below=of api] (db) {PostgreSQL};

% Monitor (acima do backend)
\node[component, above=0.9cm of api] (dash) {Supabase\\Dashboard};

\draw[arr] (gh) -- node[lbl, left] {push / PR} (ci);
\draw[arr] (ci) -- node[lbl, above] {valida} (gh);
\draw[arr] (app) -- node[lbl, above] {HTTPS} (api);
\draw[arr] (api) -- (db);
\draw[arr] (app) -- (sqlite);
\draw[arr, dashed] (dash) -- node[lbl, right] {monitora} (api);
\draw[arr, dashed] (ci) -- node[lbl, pos=0.35, above, sloped] {status CI} (gh);

\end{tikzpicture}
\caption{Arquitetura de deployment e monitoramento}
\label{fig:deployment}
\end{figure}

\subsection{Métricas e SLA}

Indicadores acompanhados na fase atual do projeto:

\begin{itemize}
    \item \textbf{Disponibilidade:} conforme plano Supabase (API e banco gerenciados).
    \item \textbf{Pipeline CI:} execução automática em \textit{push} e PR; jobs \texttt{build} e \texttt{test} obrigatórios.
    \item \textbf{Testes:} suíte Jest em \texttt{src/utils}; comando \texttt{npm run validate} localmente.
    \item \textbf{Qualidade:} \texttt{lint:ci} bloqueante no workflow; \texttt{npm run lint} ampliado para desenvolvimento.
    \item \textbf{Deploy frequency:} integração contínua via merge na \texttt{main} após CI verde.
    \item \textbf{Lead time:} revisão em PR com checks do GitHub antes do merge.
    \item \textbf{MTTR:} correção versionada, novo \textit{push} e reexecução automática do pipeline.
\end{itemize}

% Requer no preâmbulo: \usepackage{tikz}, \usetikzlibrary{positioning,arrows.meta}
%                       \usepackage{float}, \usepackage{enumitem}

\section{DevOps: Infraestrutura, Automação e Entrega Contínua}

DevOps \cite{wang2023} é uma cultura e conjunto de práticas que integra desenvolvimento de software (Dev) e operações de TI (Ops), focando em automação, colaboração e entrega contínua. No \textit{NitrusLeaf Mobile}, o repositório \texttt{Vitalliz/nitrusleaf-mobile} adota Git, GitHub Actions, backend Supabase e aplicativo Expo (React Native), com validação automática a cada alteração enviada ao repositório.

\subsection{Pipeline CI/CD}

O pipeline de integração e entrega contínua executa três \textit{jobs} no GitHub Actions (\texttt{Build}, \texttt{Test} e \texttt{Deploy}) a cada \textit{push} em qualquer branch ou em \textit{pull request} direcionado à \texttt{main}. A Figura~\ref{fig:cicd} resume o fluxo; a etapa \textbf{Monitor} representa a observação do resultado (painel do GitHub Actions e Supabase Dashboard) e o retorno de correções ao código --- ciclo de melhoria contínua.

\begin{figure}[H]
\centering
\begin{tikzpicture}[
    node distance=0.9cm and 1.1cm,
    stage/.style={
        rectangle, draw, rounded corners=3pt,
        fill=secondarycolor!20,
        text width=2.15cm, align=center,
        font=\small, minimum height=1.05cm
    },
    note/.style={font=\scriptsize, text=black!65, align=center},
    arr/.style={-{Stealth[length=2.2mm]}, thick}
]
% Linha principal
\node[stage] (code) {\textbf{Code}\\[-1pt]\scriptsize Git push / PR};
\node[stage, right=of code] (build) {\textbf{Build}\\[-1pt]\scriptsize \texttt{npm ci}};
\node[stage, right=of build] (test) {\textbf{Test}\\[-1pt]\scriptsize tipos + Jest};
\node[stage, right=of test] (deploy) {\textbf{Deploy}\\[-1pt]\scriptsize branch \texttt{main}};
\node[stage, right=of deploy] (monitor) {\textbf{Monitor}\\[-1pt]\scriptsize painéis};

\draw[arr] (code) -- (build);
\draw[arr] (build) -- (test);
\draw[arr] (test) -- (deploy);
\draw[arr] (deploy) -- (monitor);

% Detalhes abaixo (o que cada job executa de fato)
\node[note, below=0.55cm of build] (buildnote) {+\ \texttt{expo-doctor}\ (não bloqueia)};
\node[note, below=0.55cm of test] (testnote) {\texttt{typecheck},\ \texttt{lint:ci}};
\node[note, below=0.55cm of deploy] (deploynote) {integração após CI verde};

% Feedback Monitor -> Code
\draw[arr, dashed] (monitor.south) to[bend right=55]
    node[below, note, yshift=-0.15cm] {correções e novo commit} (code.south);

\end{tikzpicture}
\caption{Pipeline CI/CD implementado no repositório (\texttt{.github/workflows/ci.yml})}
\label{fig:cicd}
\end{figure}

\subsubsection{Correspondência entre etapas e jobs}

\begin{table}[H]
\centering
\caption{Etapas do pipeline e implementação no GitHub Actions}
\label{tab:cicd-jobs}
\small
\begin{tabular}{|l|l|p{7.2cm}|}
\hline
\textbf{Etapa} & \textbf{Job} & \textbf{Comandos / comportamento real} \\
\hline
Code & --- & \textit{Push} ou PR; código em \texttt{nitrusleaf-mobile/} \\
\hline
Build & \texttt{build} & \texttt{npm ci}; \texttt{npx expo-doctor} (\texttt{continue-on-error}) \\
\hline
Test & \texttt{test} & \texttt{npm run typecheck}; \texttt{npm test -- --ci}; \texttt{npm run lint:ci} \\
\hline
Deploy & \texttt{deploy} & Somente \textit{push} na \texttt{main}; confirma integração (EAS/Supabase opcionais) \\
\hline
Monitor & --- & GitHub Actions (status do workflow); Supabase Dashboard (API, Auth, BD) \\
\hline
\end{tabular}
\end{table}

O ambiente de execução utiliza \textbf{Node.js 22} no runner e a variável \texttt{FORCE\_JAVASCRIPT\_ACTIONS\_TO\_NODE24} para compatibilidade das ações oficiais do GitHub.

\subsection{Exemplo de Configuração CI/CD}

Trecho conceitual equivalente ao arquivo \texttt{.github/workflows/ci.yml}:

\begin{verbatim}
name: CI
on: [push, pull_request]

jobs:
  build:
    runs-on: ubuntu-latest
  steps:
    - uses: actions/checkout@v4
    - uses: actions/setup-node@v4
      with:
        node-version: "22"
    - run: npm ci
    - run: npx expo-doctor
      continue-on-error: true

  test:
    needs: build
  steps:
    - run: npm run typecheck
    - run: npm test -- --ci
    - run: npm run lint:ci

  deploy:
    needs: test
    if: branch == main
  steps:
    - run: echo "CI verde na main"
\end{verbatim}

\textbf{Comandos locais:}
\begin{itemize}[leftmargin=*, nosep]
    \item \texttt{npm run validate} --- verificação de tipos (\texttt{tsconfig.ci.json}) + testes Jest;
    \item \texttt{npm run lint:ci} --- ESLint na camada core (\texttt{utils}, \texttt{repositories}, \texttt{services}, \texttt{contexts});
    \item \texttt{npm run lint} --- análise completa do projeto (inclui telas; pode reportar avisos legados).
\end{itemize}

\subsection{Infraestrutura como Código}

\begin{table}[H]
\centering
\caption{Ferramentas DevOps utilizadas no NitrusLeaf Mobile}
\label{tab:devops-tools}
\begin{tabular}{|l|l|p{6cm}|}
\hline
\textbf{Categoria} & \textbf{Ferramenta} & \textbf{Finalidade} \\
\hline
Controle de Versão & Git + GitHub & \texttt{main} e branches da equipe via PR \\
CI/CD & GitHub Actions & Jobs \texttt{build}, \texttt{test} e \texttt{deploy} \\
Backend (BaaS) & Supabase & PostgreSQL, Auth JWT e API REST \\
App Mobile & Expo SDK 54 + React Native & iOS, Android e Web \\
Persistência local & \texttt{expo-sqlite} & SQLite no dispositivo (offline) \\
Testes & Jest + jest-expo & Testes em \texttt{src/utils/\_\_tests\_\_} \\
Qualidade & ESLint (\texttt{lint:ci}) & Validação automática da camada core \\
Configuração & \texttt{.env} / \texttt{.env.example} & Chaves Supabase (não versionadas) \\
\hline
\end{tabular}
\end{table}

\subsection{Infraestrutura do Sistema}

\subsubsection{Componentes por ambiente}

\begin{table}[H]
\centering
\caption{Componentes de infraestrutura}
\label{tab:infrastructure}
\begin{tabular}{|l|p{10cm}|}
\hline
\textbf{Componente} & \textbf{Descrição} \\
\hline
\textbf{Desenvolvimento local} &
\begin{itemize}[leftmargin=*, nosep]
    \item Node.js $\geq$ 20 (CI usa 22)
    \item Variáveis em \texttt{.env} conforme \texttt{.env.example}
    \item Branches dos membros integradas à \texttt{main} via pull request com CI verde
\end{itemize} \\
\hline
\textbf{Nuvem (Supabase)} &
\begin{itemize}[leftmargin=*, nosep]
    \item PostgreSQL gerenciado, autenticação e API REST
    \item HTTPS entre app e backend
    \item Painel para logs, usuários e consulta ao banco (etapa \textbf{Monitor})
\end{itemize} \\
\hline
\textbf{Dispositivos móveis} &
\begin{itemize}[leftmargin=*, nosep]
    \item Binário/servidor de desenvolvimento via Expo
    \item \texttt{expo-sqlite} para cache local
    \item Camada \texttt{src/repositories} para acesso local e remoto
\end{itemize} \\
\hline
\textbf{Segurança} &
\begin{itemize}[leftmargin=*, nosep]
    \item TLS nas requisições à API
    \item Sessão JWT via Supabase Auth
    \item Chave \textit{anon} no cliente; \texttt{.env} ignorado pelo Git
\end{itemize} \\
\hline
\end{tabular}
\end{table}

\subsubsection{Fluxo de dados}

\begin{itemize}
    \item \textbf{Usuário mobile:} autentica no Supabase Auth e consome dados via API.
    \item \textbf{Banco remoto:} PostgreSQL no Supabase.
    \item \textbf{Cache local:} SQLite (\texttt{expo-sqlite}) no aparelho.
    \item \textbf{CI/CD:} cada \textit{push} dispara build e testes; merge na \texttt{main} após pipeline bem-sucedido.
    \item \textbf{Monitoramento:} status dos workflows no GitHub; métricas e saúde da API no Supabase Dashboard.
\end{itemize}

\subsection{Arquitetura de Deploy}

A Figura~\ref{fig:deployment} separa validação do código (CI), execução do app e backend em nuvem. O GitHub Actions não hospeda o aplicativo: ele \textit{valida} o repositório antes da integração na \texttt{main}.

\begin{figure}[H]
\centering
\begin{tikzpicture}[
    node distance=1.5cm and 1.2cm,
    component/.style={
        rectangle, draw, rounded corners=3pt,
        fill=secondarycolor!20,
        text width=2.6cm, align=center,
        font=\small, minimum height=1.1cm
    },
    arr/.style={-{Stealth[length=2mm]}, thick},
    lbl/.style={font=\scriptsize, midway}
]
% Coluna CI
\node[component] (gh) {GitHub\\Repositório};
\node[component, below=of gh] (ci) {GitHub\\Actions};

% Coluna runtime
\node[component, right=3.2cm of gh] (app) {App Expo\\(dispositivo)};
\node[component, below=of app] (sqlite) {SQLite\\(\texttt{expo-sqlite})};

% Coluna backend
\node[component, right=3.2cm of app] (api) {Supabase\\Auth + API};
\node[component, below=of api] (db) {PostgreSQL};

% Monitor (acima do backend)
\node[component, above=0.9cm of api] (dash) {Supabase\\Dashboard};

\draw[arr] (gh) -- node[lbl, left] {push / PR} (ci);
\draw[arr] (ci) -- node[lbl, above] {valida} (gh);
\draw[arr] (app) -- node[lbl, above] {HTTPS} (api);
\draw[arr] (api) -- (db);
\draw[arr] (app) -- (sqlite);
\draw[arr, dashed] (dash) -- node[lbl, right] {monitora} (api);
\draw[arr, dashed] (ci) -- node[lbl, pos=0.35, above, sloped] {status CI} (gh);

\end{tikzpicture}
\caption{Arquitetura de deployment e monitoramento}
\label{fig:deployment}
\end{figure}

\subsection{Métricas e SLA}

Indicadores acompanhados na fase atual do projeto:

\begin{itemize}
    \item \textbf{Disponibilidade:} conforme plano Supabase (API e banco gerenciados).
    \item \textbf{Pipeline CI:} execução automática em \textit{push} e PR; jobs \texttt{build} e \texttt{test} obrigatórios.
    \item \textbf{Testes:} suíte Jest em \texttt{src/utils}; comando \texttt{npm run validate} localmente.
    \item \textbf{Qualidade:} \texttt{lint:ci} bloqueante no workflow; \texttt{npm run lint} ampliado para desenvolvimento.
    \item \textbf{Deploy frequency:} integração contínua via merge na \texttt{main} após CI verde.
    \item \textbf{Lead time:} revisão em PR com checks do GitHub antes do merge.
    \item \textbf{MTTR:} correção versionada, novo \textit{push} e reexecução automática do pipeline.
\end{itemize}

% Requer no preâmbulo: \usepackage{tikz}, \usetikzlibrary{positioning,arrows.meta}
%                       \usepackage{float}, \usepackage{enumitem}

\section{DevOps: Infraestrutura, Automação e Entrega Contínua}

DevOps \cite{wang2023} é uma cultura e conjunto de práticas que integra desenvolvimento de software (Dev) e operações de TI (Ops), focando em automação, colaboração e entrega contínua. No \textit{NitrusLeaf Mobile}, o repositório \texttt{Vitalliz/nitrusleaf-mobile} adota Git, GitHub Actions, backend Supabase e aplicativo Expo (React Native), com validação automática a cada alteração enviada ao repositório.

\subsection{Pipeline CI/CD}

O pipeline de integração e entrega contínua executa três \textit{jobs} no GitHub Actions (\texttt{Build}, \texttt{Test} e \texttt{Deploy}) a cada \textit{push} em qualquer branch ou em \textit{pull request} direcionado à \texttt{main}. A Figura~\ref{fig:cicd} resume o fluxo; a etapa \textbf{Monitor} representa a observação do resultado (painel do GitHub Actions e Supabase Dashboard) e o retorno de correções ao código --- ciclo de melhoria contínua.

\begin{figure}[H]
\centering
\begin{tikzpicture}[
    node distance=0.9cm and 1.1cm,
    stage/.style={
        rectangle, draw, rounded corners=3pt,
        fill=secondarycolor!20,
        text width=2.15cm, align=center,
        font=\small, minimum height=1.05cm
    },
    note/.style={font=\scriptsize, text=black!65, align=center},
    arr/.style={-{Stealth[length=2.2mm]}, thick}
]
% Linha principal
\node[stage] (code) {\textbf{Code}\\[-1pt]\scriptsize Git push / PR};
\node[stage, right=of code] (build) {\textbf{Build}\\[-1pt]\scriptsize \texttt{npm ci}};
\node[stage, right=of build] (test) {\textbf{Test}\\[-1pt]\scriptsize tipos + Jest};
\node[stage, right=of test] (deploy) {\textbf{Deploy}\\[-1pt]\scriptsize branch \texttt{main}};
\node[stage, right=of deploy] (monitor) {\textbf{Monitor}\\[-1pt]\scriptsize painéis};

\draw[arr] (code) -- (build);
\draw[arr] (build) -- (test);
\draw[arr] (test) -- (deploy);
\draw[arr] (deploy) -- (monitor);

% Detalhes abaixo (o que cada job executa de fato)
\node[note, below=0.55cm of build] (buildnote) {+\ \texttt{expo-doctor}\ (não bloqueia)};
\node[note, below=0.55cm of test] (testnote) {\texttt{typecheck},\ \texttt{lint:ci}};
\node[note, below=0.55cm of deploy] (deploynote) {integração após CI verde};

% Feedback Monitor -> Code
\draw[arr, dashed] (monitor.south) to[bend right=55]
    node[below, note, yshift=-0.15cm] {correções e novo commit} (code.south);

\end{tikzpicture}
\caption{Pipeline CI/CD implementado no repositório (\texttt{.github/workflows/ci.yml})}
\label{fig:cicd}
\end{figure}

\subsubsection{Correspondência entre etapas e jobs}

\begin{table}[H]
\centering
\caption{Etapas do pipeline e implementação no GitHub Actions}
\label{tab:cicd-jobs}
\small
\begin{tabular}{|l|l|p{7.2cm}|}
\hline
\textbf{Etapa} & \textbf{Job} & \textbf{Comandos / comportamento real} \\
\hline
Code & --- & \textit{Push} ou PR; código em \texttt{nitrusleaf-mobile/} \\
\hline
Build & \texttt{build} & \texttt{npm ci}; \texttt{npx expo-doctor} (\texttt{continue-on-error}) \\
\hline
Test & \texttt{test} & \texttt{npm run typecheck}; \texttt{npm test -- --ci}; \texttt{npm run lint:ci} \\
\hline
Deploy & \texttt{deploy} & Somente \textit{push} na \texttt{main}; confirma integração (EAS/Supabase opcionais) \\
\hline
Monitor & --- & GitHub Actions (status do workflow); Supabase Dashboard (API, Auth, BD) \\
\hline
\end{tabular}
\end{table}

O ambiente de execução utiliza \textbf{Node.js 22} no runner e a variável \texttt{FORCE\_JAVASCRIPT\_ACTIONS\_TO\_NODE24} para compatibilidade das ações oficiais do GitHub.

\subsection{Exemplo de Configuração CI/CD}

Trecho conceitual equivalente ao arquivo \texttt{.github/workflows/ci.yml}:

\begin{verbatim}
name: CI
on: [push, pull_request]

jobs:
  build:
    runs-on: ubuntu-latest
  steps:
    - uses: actions/checkout@v4
    - uses: actions/setup-node@v4
      with:
        node-version: "22"
    - run: npm ci
    - run: npx expo-doctor
      continue-on-error: true

  test:
    needs: build
  steps:
    - run: npm run typecheck
    - run: npm test -- --ci
    - run: npm run lint:ci

  deploy:
    needs: test
    if: branch == main
  steps:
    - run: echo "CI verde na main"
\end{verbatim}

\textbf{Comandos locais:}
\begin{itemize}[leftmargin=*, nosep]
    \item \texttt{npm run validate} --- verificação de tipos (\texttt{tsconfig.ci.json}) + testes Jest;
    \item \texttt{npm run lint:ci} --- ESLint na camada core (\texttt{utils}, \texttt{repositories}, \texttt{services}, \texttt{contexts});
    \item \texttt{npm run lint} --- análise completa do projeto (inclui telas; pode reportar avisos legados).
\end{itemize}

\subsection{Infraestrutura como Código}

\begin{table}[H]
\centering
\caption{Ferramentas DevOps utilizadas no NitrusLeaf Mobile}
\label{tab:devops-tools}
\begin{tabular}{|l|l|p{6cm}|}
\hline
\textbf{Categoria} & \textbf{Ferramenta} & \textbf{Finalidade} \\
\hline
Controle de Versão & Git + GitHub & \texttt{main} e branches da equipe via PR \\
CI/CD & GitHub Actions & Jobs \texttt{build}, \texttt{test} e \texttt{deploy} \\
Backend (BaaS) & Supabase & PostgreSQL, Auth JWT e API REST \\
App Mobile & Expo SDK 54 + React Native & iOS, Android e Web \\
Persistência local & \texttt{expo-sqlite} & SQLite no dispositivo (offline) \\
Testes & Jest + jest-expo & Testes em \texttt{src/utils/\_\_tests\_\_} \\
Qualidade & ESLint (\texttt{lint:ci}) & Validação automática da camada core \\
Configuração & \texttt{.env} / \texttt{.env.example} & Chaves Supabase (não versionadas) \\
\hline
\end{tabular}
\end{table}

\subsection{Infraestrutura do Sistema}

\subsubsection{Componentes por ambiente}

\begin{table}[H]
\centering
\caption{Componentes de infraestrutura}
\label{tab:infrastructure}
\begin{tabular}{|l|p{10cm}|}
\hline
\textbf{Componente} & \textbf{Descrição} \\
\hline
\textbf{Desenvolvimento local} &
\begin{itemize}[leftmargin=*, nosep]
    \item Node.js $\geq$ 20 (CI usa 22)
    \item Variáveis em \texttt{.env} conforme \texttt{.env.example}
    \item Branches dos membros integradas à \texttt{main} via pull request com CI verde
\end{itemize} \\
\hline
\textbf{Nuvem (Supabase)} &
\begin{itemize}[leftmargin=*, nosep]
    \item PostgreSQL gerenciado, autenticação e API REST
    \item HTTPS entre app e backend
    \item Painel para logs, usuários e consulta ao banco (etapa \textbf{Monitor})
\end{itemize} \\
\hline
\textbf{Dispositivos móveis} &
\begin{itemize}[leftmargin=*, nosep]
    \item Binário/servidor de desenvolvimento via Expo
    \item \texttt{expo-sqlite} para cache local
    \item Camada \texttt{src/repositories} para acesso local e remoto
\end{itemize} \\
\hline
\textbf{Segurança} &
\begin{itemize}[leftmargin=*, nosep]
    \item TLS nas requisições à API
    \item Sessão JWT via Supabase Auth
    \item Chave \textit{anon} no cliente; \texttt{.env} ignorado pelo Git
\end{itemize} \\
\hline
\end{tabular}
\end{table}

\subsubsection{Fluxo de dados}

\begin{itemize}
    \item \textbf{Usuário mobile:} autentica no Supabase Auth e consome dados via API.
    \item \textbf{Banco remoto:} PostgreSQL no Supabase.
    \item \textbf{Cache local:} SQLite (\texttt{expo-sqlite}) no aparelho.
    \item \textbf{CI/CD:} cada \textit{push} dispara build e testes; merge na \texttt{main} após pipeline bem-sucedido.
    \item \textbf{Monitoramento:} status dos workflows no GitHub; métricas e saúde da API no Supabase Dashboard.
\end{itemize}

\subsection{Arquitetura de Deploy}

A Figura~\ref{fig:deployment} separa validação do código (CI), execução do app e backend em nuvem. O GitHub Actions não hospeda o aplicativo: ele \textit{valida} o repositório antes da integração na \texttt{main}.

\begin{figure}[H]
\centering
\begin{tikzpicture}[
    node distance=1.5cm and 1.2cm,
    component/.style={
        rectangle, draw, rounded corners=3pt,
        fill=secondarycolor!20,
        text width=2.6cm, align=center,
        font=\small, minimum height=1.1cm
    },
    arr/.style={-{Stealth[length=2mm]}, thick},
    lbl/.style={font=\scriptsize, midway}
]
% Coluna CI
\node[component] (gh) {GitHub\\Repositório};
\node[component, below=of gh] (ci) {GitHub\\Actions};

% Coluna runtime
\node[component, right=3.2cm of gh] (app) {App Expo\\(dispositivo)};
\node[component, below=of app] (sqlite) {SQLite\\(\texttt{expo-sqlite})};

% Coluna backend
\node[component, right=3.2cm of app] (api) {Supabase\\Auth + API};
\node[component, below=of api] (db) {PostgreSQL};

% Monitor (acima do backend)
\node[component, above=0.9cm of api] (dash) {Supabase\\Dashboard};

\draw[arr] (gh) -- node[lbl, left] {push / PR} (ci);
\draw[arr] (ci) -- node[lbl, above] {valida} (gh);
\draw[arr] (app) -- node[lbl, above] {HTTPS} (api);
\draw[arr] (api) -- (db);
\draw[arr] (app) -- (sqlite);
\draw[arr, dashed] (dash) -- node[lbl, right] {monitora} (api);
\draw[arr, dashed] (ci) -- node[lbl, pos=0.35, above, sloped] {status CI} (gh);

\end{tikzpicture}
\caption{Arquitetura de deployment e monitoramento}
\label{fig:deployment}
\end{figure}

\subsection{Métricas e SLA}

Indicadores acompanhados na fase atual do projeto:

\begin{itemize}
    \item \textbf{Disponibilidade:} conforme plano Supabase (API e banco gerenciados).
    \item \textbf{Pipeline CI:} execução automática em \textit{push} e PR; jobs \texttt{build} e \texttt{test} obrigatórios.
    \item \textbf{Testes:} suíte Jest em \texttt{src/utils}; comando \texttt{npm run validate} localmente.
    \item \textbf{Qualidade:} \texttt{lint:ci} bloqueante no workflow; \texttt{npm run lint} ampliado para desenvolvimento.
    \item \textbf{Deploy frequency:} integração contínua via merge na \texttt{main} após CI verde.
    \item \textbf{Lead time:} revisão em PR com checks do GitHub antes do merge.
    \item \textbf{MTTR:} correção versionada, novo \textit{push} e reexecução automática do pipeline.
\end{itemize}

% Requer no preâmbulo: \usepackage{tikz}, \usetikzlibrary{positioning,arrows.meta}
%                       \usepackage{float}, \usepackage{enumitem}

\section{DevOps: Infraestrutura, Automação e Entrega Contínua}

DevOps \cite{wang2023} é uma cultura e conjunto de práticas que integra desenvolvimento de software (Dev) e operações de TI (Ops), focando em automação, colaboração e entrega contínua. No \textit{NitrusLeaf Mobile}, o repositório \texttt{Vitalliz/nitrusleaf-mobile} adota Git, GitHub Actions, backend Supabase e aplicativo Expo (React Native), com validação automática a cada alteração enviada ao repositório.

\subsection{Pipeline CI/CD}

O pipeline de integração e entrega contínua executa três \textit{jobs} no GitHub Actions (\texttt{Build}, \texttt{Test} e \texttt{Deploy}) a cada \textit{push} em qualquer branch ou em \textit{pull request} direcionado à \texttt{main}. A Figura~\ref{fig:cicd} resume o fluxo; a etapa \textbf{Monitor} representa a observação do resultado (painel do GitHub Actions e Supabase Dashboard) e o retorno de correções ao código --- ciclo de melhoria contínua.

\begin{figure}[H]
\centering
\begin{tikzpicture}[
    node distance=0.9cm and 1.1cm,
    stage/.style={
        rectangle, draw, rounded corners=3pt,
        fill=secondarycolor!20,
        text width=2.15cm, align=center,
        font=\small, minimum height=1.05cm
    },
    note/.style={font=\scriptsize, text=black!65, align=center},
    arr/.style={-{Stealth[length=2.2mm]}, thick}
]
% Linha principal
\node[stage] (code) {\textbf{Code}\\[-1pt]\scriptsize Git push / PR};
\node[stage, right=of code] (build) {\textbf{Build}\\[-1pt]\scriptsize \texttt{npm ci}};
\node[stage, right=of build] (test) {\textbf{Test}\\[-1pt]\scriptsize tipos + Jest};
\node[stage, right=of test] (deploy) {\textbf{Deploy}\\[-1pt]\scriptsize branch \texttt{main}};
\node[stage, right=of deploy] (monitor) {\textbf{Monitor}\\[-1pt]\scriptsize painéis};

\draw[arr] (code) -- (build);
\draw[arr] (build) -- (test);
\draw[arr] (test) -- (deploy);
\draw[arr] (deploy) -- (monitor);

% Detalhes abaixo (o que cada job executa de fato)
\node[note, below=0.55cm of build] (buildnote) {+\ \texttt{expo-doctor}\ (não bloqueia)};
\node[note, below=0.55cm of test] (testnote) {\texttt{typecheck},\ \texttt{lint:ci}};
\node[note, below=0.55cm of deploy] (deploynote) {integração após CI verde};

% Feedback Monitor -> Code
\draw[arr, dashed] (monitor.south) to[bend right=55]
    node[below, note, yshift=-0.15cm] {correções e novo commit} (code.south);

\end{tikzpicture}
\caption{Pipeline CI/CD implementado no repositório (\texttt{.github/workflows/ci.yml})}
\label{fig:cicd}
\end{figure}

\subsubsection{Correspondência entre etapas e jobs}

\begin{table}[H]
\centering
\caption{Etapas do pipeline e implementação no GitHub Actions}
\label{tab:cicd-jobs}
\small
\begin{tabular}{|l|l|p{7.2cm}|}
\hline
\textbf{Etapa} & \textbf{Job} & \textbf{Comandos / comportamento real} \\
\hline
Code & --- & \textit{Push} ou PR; código em \texttt{nitrusleaf-mobile/} \\
\hline
Build & \texttt{build} & \texttt{npm ci}; \texttt{npx expo-doctor} (\texttt{continue-on-error}) \\
\hline
Test & \texttt{test} & \texttt{npm run typecheck}; \texttt{npm test -- --ci}; \texttt{npm run lint:ci} \\
\hline
Deploy & \texttt{deploy} & Somente \textit{push} na \texttt{main}; confirma integração (EAS/Supabase opcionais) \\
\hline
Monitor & --- & GitHub Actions (status do workflow); Supabase Dashboard (API, Auth, BD) \\
\hline
\end{tabular}
\end{table}

O ambiente de execução utiliza \textbf{Node.js 22} no runner e a variável \texttt{FORCE\_JAVASCRIPT\_ACTIONS\_TO\_NODE24} para compatibilidade das ações oficiais do GitHub.

\subsection{Exemplo de Configuração CI/CD}

Trecho conceitual equivalente ao arquivo \texttt{.github/workflows/ci.yml}:

\begin{verbatim}
name: CI
on: [push, pull_request]

jobs:
  build:
    runs-on: ubuntu-latest
  steps:
    - uses: actions/checkout@v4
    - uses: actions/setup-node@v4
      with:
        node-version: "22"
    - run: npm ci
    - run: npx expo-doctor
      continue-on-error: true

  test:
    needs: build
  steps:
    - run: npm run typecheck
    - run: npm test -- --ci
    - run: npm run lint:ci

  deploy:
    needs: test
    if: branch == main
  steps:
    - run: echo "CI verde na main"
\end{verbatim}

\textbf{Comandos locais:}
\begin{itemize}[leftmargin=*, nosep]
    \item \texttt{npm run validate} --- verificação de tipos (\texttt{tsconfig.ci.json}) + testes Jest;
    \item \texttt{npm run lint:ci} --- ESLint na camada core (\texttt{utils}, \texttt{repositories}, \texttt{services}, \texttt{contexts});
    \item \texttt{npm run lint} --- análise completa do projeto (inclui telas; pode reportar avisos legados).
\end{itemize}

\subsection{Infraestrutura como Código}

\begin{table}[H]
\centering
\caption{Ferramentas DevOps utilizadas no NitrusLeaf Mobile}
\label{tab:devops-tools}
\begin{tabular}{|l|l|p{6cm}|}
\hline
\textbf{Categoria} & \textbf{Ferramenta} & \textbf{Finalidade} \\
\hline
Controle de Versão & Git + GitHub & \texttt{main} e branches da equipe via PR \\
CI/CD & GitHub Actions & Jobs \texttt{build}, \texttt{test} e \texttt{deploy} \\
Backend (BaaS) & Supabase & PostgreSQL, Auth JWT e API REST \\
App Mobile & Expo SDK 54 + React Native & iOS, Android e Web \\
Persistência local & \texttt{expo-sqlite} & SQLite no dispositivo (offline) \\
Testes & Jest + jest-expo & Testes em \texttt{src/utils/\_\_tests\_\_} \\
Qualidade & ESLint (\texttt{lint:ci}) & Validação automática da camada core \\
Configuração & \texttt{.env} / \texttt{.env.example} & Chaves Supabase (não versionadas) \\
\hline
\end{tabular}
\end{table}

\subsection{Infraestrutura do Sistema}

\subsubsection{Componentes por ambiente}

\begin{table}[H]
\centering
\caption{Componentes de infraestrutura}
\label{tab:infrastructure}
\begin{tabular}{|l|p{10cm}|}
\hline
\textbf{Componente} & \textbf{Descrição} \\
\hline
\textbf{Desenvolvimento local} &
\begin{itemize}[leftmargin=*, nosep]
    \item Node.js $\geq$ 20 (CI usa 22)
    \item Variáveis em \texttt{.env} conforme \texttt{.env.example}
    \item Branches dos membros integradas à \texttt{main} via pull request com CI verde
\end{itemize} \\
\hline
\textbf{Nuvem (Supabase)} &
\begin{itemize}[leftmargin=*, nosep]
    \item PostgreSQL gerenciado, autenticação e API REST
    \item HTTPS entre app e backend
    \item Painel para logs, usuários e consulta ao banco (etapa \textbf{Monitor})
\end{itemize} \\
\hline
\textbf{Dispositivos móveis} &
\begin{itemize}[leftmargin=*, nosep]
    \item Binário/servidor de desenvolvimento via Expo
    \item \texttt{expo-sqlite} para cache local
    \item Camada \texttt{src/repositories} para acesso local e remoto
\end{itemize} \\
\hline
\textbf{Segurança} &
\begin{itemize}[leftmargin=*, nosep]
    \item TLS nas requisições à API
    \item Sessão JWT via Supabase Auth
    \item Chave \textit{anon} no cliente; \texttt{.env} ignorado pelo Git
\end{itemize} \\
\hline
\end{tabular}
\end{table}

\subsubsection{Fluxo de dados}

\begin{itemize}
    \item \textbf{Usuário mobile:} autentica no Supabase Auth e consome dados via API.
    \item \textbf{Banco remoto:} PostgreSQL no Supabase.
    \item \textbf{Cache local:} SQLite (\texttt{expo-sqlite}) no aparelho.
    \item \textbf{CI/CD:} cada \textit{push} dispara build e testes; merge na \texttt{main} após pipeline bem-sucedido.
    \item \textbf{Monitoramento:} status dos workflows no GitHub; métricas e saúde da API no Supabase Dashboard.
\end{itemize}

\subsection{Arquitetura de Deploy}

A Figura~\ref{fig:deployment} separa validação do código (CI), execução do app e backend em nuvem. O GitHub Actions não hospeda o aplicativo: ele \textit{valida} o repositório antes da integração na \texttt{main}.

\begin{figure}[H]
\centering
\begin{tikzpicture}[
    node distance=1.5cm and 1.2cm,
    component/.style={
        rectangle, draw, rounded corners=3pt,
        fill=secondarycolor!20,
        text width=2.6cm, align=center,
        font=\small, minimum height=1.1cm
    },
    arr/.style={-{Stealth[length=2mm]}, thick},
    lbl/.style={font=\scriptsize, midway}
]
% Coluna CI
\node[component] (gh) {GitHub\\Repositório};
\node[component, below=of gh] (ci) {GitHub\\Actions};

% Coluna runtime
\node[component, right=3.2cm of gh] (app) {App Expo\\(dispositivo)};
\node[component, below=of app] (sqlite) {SQLite\\(\texttt{expo-sqlite})};

% Coluna backend
\node[component, right=3.2cm of app] (api) {Supabase\\Auth + API};
\node[component, below=of api] (db) {PostgreSQL};

% Monitor (acima do backend)
\node[component, above=0.9cm of api] (dash) {Supabase\\Dashboard};

\draw[arr] (gh) -- node[lbl, left] {push / PR} (ci);
\draw[arr] (ci) -- node[lbl, above] {valida} (gh);
\draw[arr] (app) -- node[lbl, above] {HTTPS} (api);
\draw[arr] (api) -- (db);
\draw[arr] (app) -- (sqlite);
\draw[arr, dashed] (dash) -- node[lbl, right] {monitora} (api);
\draw[arr, dashed] (ci) -- node[lbl, pos=0.35, above, sloped] {status CI} (gh);

\end{tikzpicture}
\caption{Arquitetura de deployment e monitoramento}
\label{fig:deployment}
\end{figure}

\subsection{Métricas e SLA}

Indicadores acompanhados na fase atual do projeto:

\begin{itemize}
    \item \textbf{Disponibilidade:} conforme plano Supabase (API e banco gerenciados).
    \item \textbf{Pipeline CI:} execução automática em \textit{push} e PR; jobs \texttt{build} e \texttt{test} obrigatórios.
    \item \textbf{Testes:} suíte Jest em \texttt{src/utils}; comando \texttt{npm run validate} localmente.
    \item \textbf{Qualidade:} \texttt{lint:ci} bloqueante no workflow; \texttt{npm run lint} ampliado para desenvolvimento.
    \item \textbf{Deploy frequency:} integração contínua via merge na \texttt{main} após CI verde.
    \item \textbf{Lead time:} revisão em PR com checks do GitHub antes do merge.
    \item \textbf{MTTR:} correção versionada, novo \textit{push} e reexecução automática do pipeline.
\end{itemize}
